\section{The Others}

\begin{quotex}
The deceased are around us and participate in our lives. Even our thoughts, desires, insights are influenced by them. We have to be aware of this. Indeed, they move and live around us! During our sleep we don't sense the physical objects that are beside us, so during our wake-up life we don't feel the dead around us. Only the state of our consciousness separates us from them. Dialogue with the dead is a very concrete activity and no expedient can replace human efforts to transform pain into gratitude. Gratitude for all that has been received and shared in life.
\flright{\textsc{Rudolf Steiner}}
\end{quotex}


\emph{The Others}, starring \textbf{Nicole Kidman}, is one of the creepiest horror movies ever produced. It does not rely on blood, gore, or sudden surprise moves; rather, the psychodrama of an evil event persists beyond the grave. With All Soul's Day coming up, it is appropriate to explore the various postmortem states of the human soul. This will be done from the metaphysical perspective and only tangentially related to the exoteric and theological dogmas.

\paragraph{Kamaloka}


It is not uncommon for the souls of the dead to be associated with the location of their death or burial. In the Ancient City\footnote{\url{https://gornahoor.net/?p=2507}}, the dead were buried with utensils, clothing, and weapons.

\textbf{Valentin Tomberg} describes this state:


\begin{quotex}
Kamaloka means ``desire world,'' a partially material realm that is usually invisible to human beings. It surrounds and encloses our physical Earth and is the dwelling place of the astral forms of dead human and other beings. It is in Kamaloka that the second death takes place, after which the liberated part of the human being enters Devachan.
\end{quotex}


Every death is a birth into another state. \textbf{Rene Guenon}, in \emph{Man and his Becoming}, asserts:


\begin{quotex}
It must be realized that the idea of death is essentially synonymous with a change of state, which is its widest acceptation; and when it is said that the being has virtually attained immortality, that is taken to mean that it will not need to pass through further conditioned states different from the human state, or to traverse other cycles of manifestation.
\end{quotex}


In the \emph{Symbolism of the Cross}, there is a chapter on the metaphysical equivalence of birth and death.

\begin{quotex}
These two phenomena accordingly accompany and complete each other: human birth is the immediate result of a death (to another state); human death is the immediate cause of a birth (likewise into another state) . Neither of these circumstances can ever occur without the other. And, as time does not exist here, it can be affirmed that, between the intrinsic value of the phenomenon birth and the intrinsic value of the phenomenon death, there is metaphysical identity.
\end{quotex}


Devachan is a more pleasant state of being, but there are other possibilities. The Traditional teaching of the West\footnote{\url{https://gornahoor.net/?p=7568}}, shown for example by \textbf{Maximos the Confessor}, is that the postmortem options refer to states of being, not to actual places. This is not commonly acknowledged in our time. Hell, Limbo, Paradise, and the Empyrean represent various states. What is often forgotten is that the states are temporary, since they are way stations prior to Resurrection.

\paragraph{Hell}


There are different levels of hell, each appropriate to the soul's state of being. In hell, suffering is not alleviated. There is no second death in hell.

\paragraph{Limbo}


This is the state of natural happiness, in which normal desires are satisfied but there is no supernatural happiness nor awareness of the presence of God. This is appropriate to those who have been equally balanced between good and evil. For example, philosophies like Stoicism or Epicureanism promote such a balance. In the Ancient City, great men and heroes will find their destiny in Limbo. They may have been wise rulers, but their greatness was not related to a life of virtue.

\paragraph{Earthly Paradise}


Paradise is a return to the Primordial State, the status quo ante before the Fall. As such, it is prolongation of the human state into more subtle realms. This is what most people understand by ``heaven''. This state is called Salvation.

Guenon offers this explanation:

\begin{quotationx}
Westerners merely mean an extension of the possibilities of the human order, consisting in an indefinite prolongation of life (what the Far-Eastern tradition calls `longevity') under conditions which are to a certain degree transposed but which always remain more or less similar to those of terrestrial existence, since they likewise concern the human individuality.

`Salvation' in the religious sense given to the word by Western people, being the fruit of certain actions. This accords perfectly with the western conception of immortality, which is simply an indefinite prolongation of individual life transposed into the subtle order.
\end{quotationx}


In this state, desires are satisfied and the individual will interact with those he had known on earth.

\paragraph{Empyrean}


This is also known as the Heavenly Jerusalem. Beyond the Earthly Paradise, there are the spheres of the planets and the stars, representing more subtle states of the being. Ultimately, in the Empyrean, the Being is beyond all individual conditions, including time and space. There one achieves the Beatific Vision. It is possible for some saints to achieve the Beatific Vision directly from the human state.

\paragraph{Degrees of Existence}


It must not be assumed that there are always discrete jumps. In \emph{The Symbolism of the Cross}, Guenon describes how the degrees are actually infinitesimally close to each other. So Dante and Beatrice pass through all these states on the path to the Heavenly Jerusalem.


\flrightit{Posted on 2021-10-27 by Cologero}

\begin{center}* * *\end{center}

\begin{footnotesize}
\begin{sffamily}
\texttt{Jacek on 2021-10-28 at 11:01 said:}

Did Guenon subscribe to the idea of ``second death''?

Where does ``second death'' fit into all of this?

\hfill

\texttt{Reza on 2021-10-28 at 11:09 said:}

``There is no second death in hell''

There is a contradiction here. If we are ``born'' into hell, then necessarily we must die to it.

\hfill

\texttt{Cologero on 2021-10-28 at 21:53 said:}

There is no contradiction. Hell is a prolongation of the human state which may be immortal. Guenon explains:

\begin{quote}\small
We are especially alluding to the extension of the idea of life which is implied in the point of view of the Western religions, and which in fact relates to possibilities contained in a prolongation of human individuality; as we have explained elsewhere, this is what the Far-Eastern tradition refers to under the name of `longevity'.
\end{quote}


That prolongation persists:

\begin{quote}\small
By the corresponding word Westerners merely mean an extension of the possibilities of the human order, consisting in an indefinite prolongation of life under conditions which are to a certain degree transposed but which always remain more or less similar to those of terrestrial existence, since they likewise concern the human individuality. Now in the present instance the state described is still an individual state and nevertheless it is said that immortality can be obtained therein.
\end{quote}


Guenon even explains why that is not a contradiction:

\begin{quote}\small
This may appear \textbf{inconsistent} with what we have just remarked, since it might be supposed that relative immortality only is meant, understood according to the Western sense: actually however \textbf{that is not the case}.
\end{quote}


Death means the transition from one state of being to another, provided the being has the effective knowledge to make the transition. It is difficult to see how the being in hell, which is a static state, can possibly grow in knowledge.

\begin{quote}\small
The transition from each stage to the next only becoming possible for the being who has obtained the corresponding degree of effective knowledge.
\end{quote}


\hfill

\texttt{Cologero on 2021-10-28 at 21:58 said:}

Guenon does not use the term ``second death'' exactly. However, in the \emph{Symbolism of the Cross}, he points out the metaphysical equivalence of birth and death. In other words, death, as the end of one state, is actually the birth into another state. That is what Tomberg is alluding to: The first human death is the transition to the intermediate state (roughly equivalent to purgatory). The being dies to that state and enters a higher state.

\hfill

\texttt{Paulo Adolpho Zeymer Netto on 2021-10-31 at 18:56 said:}

In Evola's book ``Introduction to Magic'', he starts from a view similar to that of George Gurdjieff, who says that human beings do not have an immortal soul, but must acquire one in life, the creation of the astral body. What does Guenon have to say about this?

\hfill

\end{sffamily}
\end{footnotesize}

